%!TEX root = ../thesis.tex
%*******************************************************************************
%*********************************** First Chapter *****************************
%*******************************************************************************

\ifpdf
    \graphicspath{{Chapter1/Figs/Raster/}{Chapter1/Figs/PDF/}{Chapter1/Figs/}}
\else
    \graphicspath{{Chapter1/Figs/Vector/}{Chapter1/Figs/}}
\fi


%********************************** %First Section  **************************************
\chapter{Introduction}
% \nomenclature[z-PPC]{PPC}{Particles per cell}

Computed tomography (CT) is widely used in medicine and industry to investigate
internal structure non-destructively. X-rays are directed through an object from
multiple angles and their attenuation is measured at a detector. In medical
imaging, reducing the number of views or photon flux can lower patient dose, but
also makes reconstruction more ill-posed and increases the influence of sampling
artefacts and noise \cite{sidky_accurate_2006, chang_modeling_2017}. Industrial
systems and research applications can often use more views and higher energies,
yet acquisition time, motion, restricted access and radiation-sensitive
specimens create related constraints \cite{sun_realisation_2023}.

Learned image priors can regularise these underdetermined problems. Diffusion models are particularly expressive, but conventional whole-image models require substantial memory and training data. Hu et al. proposed the Patch Diffusion Inverse Solver (PaDIS), which learns position-conditioned patch scores and combines randomly shifted patch partitions during sampling \cite{hu_learning_2024}. Hu et al. report that PaDIS reduces training cost and outperforms an otherwise comparable whole-image diffusion prior on sparse-view CT. These claims matter for high-resolution medical imaging with small datasets, but the original PaDIS paper and its released implementation leave several experimental choices ambiguous or inconsistent.

The primary objectives are to:
\begin{itemize}[leftmargin=1.5em,itemsep=0.1em,topsep=0.2em,parsep=0pt]
    \item Reproduce the core PaDIS result on a different public CT dataset using a physically specified fan-beam operator.
    \item Compare the reproductions produced by a re-implementation of those methods desribed in the original PaDIS paper and codebase and evaluate the operator-normalised extension introduced here.
    \item Quantify the effects of patch size, training-set size, positional encoding, initialisation and sampler choice.
\end{itemize}
To address these objectives, we implemented the complete pipeline in LION \cite{kiss_benchmarking_2025,noauthor_cambridgecialion_2026} using patient-separated LIDC-IDRI splits \cite{armato_iii_data_2015-1,armato_lung_2011}. The evaluation includes FDK, TV and plug-and-play baselines, alternative diffusion samplers, fixed patch-assembly methods, PaDIS and whole-image diffusion. Experiments cover 8, 20 and 60 views, limited-angle acquisition, and native $512\times512$ reconstruction. Quantitative results include image-level bootstrap uncertainty and measurement-consistency metrics.

This work therefore tests the generalisability of the central claims made by Hu et al. \cite{hu_learning_2024} across different data and geometry rather than seeking numerical identity. The remaining chapters present the background, methods, results and discussion before the conclusion.
